
\chapter*{General Introduction}
\markboth{General Introduction}{} %pour afficher l'entete
\addcontentsline{toc}{chapter}{General Introduction }

Over the past decades, the growing demand for safe, portable, and cost-effective medical technologies has driven the emergence of innovative biomedical imaging modalities. These solutions aim to provide real-time monitoring and diagnosis while overcoming the limitations of conventional imaging techniques such as CT, MRI, or ultrasound, which may be expensive, bulky, or potentially harmful in repeated use. Among the promising directions, electrical-based imaging approaches are attracting significant attention, as they combine simplicity of instrumentation with the ability to capture functional physiological information.

Within this context, Electrical Impedance Tomography (EIT) has gained interest for a wide class of biomedical applications, including pulmonary monitoring, cancer detection, and functional imaging. However, most conventional EIT systems rely on single-frequency (monofrequency) excitation, which may limit their capacity to distinguish between tissues with similar impedance values at one frequency. This limitation directly affects the sensitivity and specificity of the imaging, reducing diagnostic reliability in clinical applications where fine tissue differentiation is crucial. To address this challenge, multi-frequency excitation has emerged as a promising approach. By measuring impedance across multiple frequencies, it becomes possible to exploit the frequency-dependent electrical properties of biological tissues, thereby improving image quality and diagnostic performance.

Nevertheless, the design and implementation of multi-frequency EIT systems introduce a number of engineering challenges. These include the generation of stable multi-tone signals, robust data acquisition strategies, advanced signal processing techniques to extract relevant features, and efficient image reconstruction algorithms. This graduation project is dedicated to the development of a multi-frequency EIT system tailored for medical imaging. The objective is to design and implement a hardware platform capable of generating and acquiring multi-frequency signals, establish a dynamic signal processing pipeline, and experimentally evaluate system performance against conventional monofrequency approaches. The work spans the entire engineering process, from system modeling and circuit design to PCB realization, data processing, and performance analysis, all while considering biomedical instrumentation constraints such as safety, signal integrity, and measurement accuracy.

The remainder of this report is organized as follows: Chapter one presents the host institution, the project context, and the objectives; Chapter two reviews the theoretical foundations of EIT, its biomedical applications, and the motivation for multi-frequency systems; Chapter three describes the proposed system architecture and its implementation, from modeling to hardware realization; Chapter four discusses the signal processing methods and experimental results, followed by an analysis of limitations and possible improvements. Through this project, we aim to contribute to the development of accessible and effective biomedical imaging solutions while reinforcing the practical engineering expertise required for designing complex instrumentation systems.

